\begin{figure}[t]
\centering
\fbox{\begin{minipage}{0.95\columnwidth}
\centering
\textbf{Start: marine component candidate for AM-HEA}\\[2pt]
$\downarrow$\\[-2pt]
Complex internal geometry (conformal channels, integrated manifolds)?\\
\textbf{yes} $\rightarrow$ LPBF HEA (Section~\ref{sec:am}); \textbf{no} $\downarrow$\\[2pt]
Large/near-net-shape part to be machined to finish?\\
\textbf{yes} $\rightarrow$ WAAM HEA + machining (wire over-alloying per Section~\ref{sec:a13}); \textbf{no} $\downarrow$\\[2pt]
Repair, upgrade, or hybrid economy (existing metallic part)?\\
\textbf{yes} $\rightarrow$ DED/LMD HEA cladding or graded deposit on steel/bronze substrate; \textbf{no} $\downarrow$\\[2pt]
Small batch, no qualification data, non-critical service?\\
\textbf{yes} $\rightarrow$ retain incumbent (NiAl bronze / 316L); re-evaluate when exposure data exist\\[2pt]
$\downarrow$ (all AM routes)\\[-2pt]
Corrosion evidence gate: G52 exposure matrix (Table~\ref{tab:exposure}) + crevice/G192 screening passed?\\
\textbf{no} $\rightarrow$ laboratory evidence first; \textbf{yes} $\rightarrow$ class submission (Section~\ref{sec:applications})\\[2pt]
$\downarrow$\\[-2pt]
Life-cycle check: interval extension and critical-element burden justify substitution? \cite{nass2025marine}
\end{minipage}}
\caption{Decision tree for route selection and deployment gating of AM-HEA marine components, linked to the indicative cost model (Table~\ref{tab:costmodel}). All decision logic is proposal \prop; the cost anchors behind it are flagged in Section~\ref{sec:a14}.}
\label{fig:decision_tree}
\end{figure}
